\def\level{Première}  % Classe à droite
\def\course{NSI}
\def\chaptername{Adresses IP} % Titre au centre
\def\docname{ARSE05}        % Identifiant à gauche

\pagestyle{cours_FP}
\makeheader{} % Affiche le bandeau


\begin{definition}[Réseau informatique]{}
	C'est un ensemble de machines (ordinateurs, smart phones, imprimantes, box, objets connectés \dots) reliés entre eux. Chaque objet possède une \textbf{carte réseau} et est connecté au réseau (avec ou sans fil).
\end{definition}



Pour qu'un équipement puisse être connecté à un réseau, il lui fait un interface réseau. Cet interface est soit filaire via une connexion de type RJ45 ou sans fil (wifi ou bluetooth).

\begin{exemple}[Connexion filaire]{}
\begin{minipage}{0.40\linewidth}
	\begin{itemize}
		\item une carte réseau dans un PC fixe
		\item un adaptateur USB-RJ45
		\item un prise RJ45
	\end{itemize}
\end{minipage}\hfill
\begin{minipage}{0.15\linewidth}
	\includegraphics[width = 0.9\linewidth]{cr1.jpg}
\end{minipage}\hfill
\begin{minipage}{0.15\linewidth}
	\includegraphics[width = 0.9\linewidth]{cr2.jpg}
\end{minipage}\hfill
\begin{minipage}{0.15\linewidth}
	\includegraphics[width = 0.9\linewidth]{cr3.jpg}
\end{minipage}
\end{exemple}

\begin{easybox}[Une carte réseau -- une adresse physique]{}
Chaque interface réseau doit posséder une \textbf{adresse physique}, nommée adresse \verb+MAC+: elle est liée à la carte réseau et identifie de manière unique celle-ci. Cette adresse ne peut pas être modifiée (facilement).

C'est une suite de six nombres en hexadécimal. par exemple: \verb+94-BB-43-42-A3-A5+.
\end{easybox}


\begin{easybox}[Constitution d'un réseau]{}

	\begin{minipage}{0.48\linewidth}
		\begin{itemize}
	\item un \textbf{routeur}: la porte d'entrée du réseau.
	\item un \textbf{commutateur} (switch): ``multiprise'' 
	
	pour connecter plusieurs objets au routeur
	\item des \textbf{postes clients}: les PC des utilisateurs
	\item un \textbf{serveur}: permet de stocker des données accessibles par tout le monde
	\item une \textbf{borne wifi} créer un nouveau sous-réseau
\end{itemize}

	\end{minipage}\hfill{}
	\begin{minipage}{0.48\linewidth}
\includegraphics[width=1\linewidth]{reseau2.jpg}

	\end{minipage}


	
\end{easybox}


\begin{center}
	\includegraphics[width = 0.9\linewidth]{ScreenShot_20240601232651.jpeg}

\end{center}
